\section{Komplexe Analysis}

\subsection{Ausdruck für die komplex konjugierte Geschwindigkeit $\overline{v(z)}$}

Um den Ausdruck~\eqref{joukowski:auftrieb:auftrieb} weiter aufzulösen, muss die oben verwendete Funktion $v(z)$ genauer untersucht werden.
Einige Eigenschaften von $v(z)$ helfen uns dabei einen passenden Ausdruck dafür aufzustellen.
Wir haben vorausgesetz, dass im Unendlichen die ungestörte Anströmungsgeschwindigkeit $v(z) = v_\infty \in \mathbb{R}$ ergibt.
Zudem muss es laut Beispiel~\ref{buch:integration:wege:bsp:kreisintegral} eine einfache Polstelle im inneren des Flügelprofils haben,
da sonst das Wegintegral um das Flügelprofil konstant $0$ wäre,
was darauf hindeuten würde, dass keinen Auftrieb erzeilt werden kann.
Diese beiden Eigenschaften der Geschwindigkeit $v(z)$ führen dazu, 
dass die komplex konjugierte Geschwidigkeit $\overline{v(z)}$ als Laurent-Reihe der Form

\[
	\overline{v(z)} 
	= 
	a_0 + \frac{a_1}{z} + \frac{a_2}{z^2} + \frac{a_3}{z^3} + \dots 
	= 
	\sum_{k=0}^{\infty} \frac{a_k}{z^k}
\]
geschrieben werden kann.

Da die Geschwindigkeit $v(z)$ überall ausser im Punkt $z_0 = 0$ definiert ist, muss auch der Reihenausdruck diese Eigenschaft haben.
Dies wird zum Beispiel von einer Taylor-Reihe nicht erfüllt, da es damit unmöglich ist eine Definitionslücke bei $z_0 = 0$ zu erreichen,
weshalb die Laurent-Reihe gewählt wurde. Auch der konstante Wert im Unendlichen bereitet der Taylor-Reihe schwierigkeiten.

Im Allgemeinen sind die Koeffizienten $a_k \in \mathbb{C}$, einzelne beschränken sich jedoch auf $a_k \in \mathbb{R}$, wie $a_0$.
Der Koeffizient $a_0$ entspricht direkt der Geschwindigkeit im Unendlichkeit, weshalb er besser als $v_\infty$ geschrieben werden soll.
Für die komplex konjugierte Geschwindigkeit $\overline{v(z)}$ ergibt sich somit die besser formulierte Reihenentwicklung

\begin{equation}
	\overline{v(z)}
	=
	v_\infty + \sum_{k=1}^{\infty} \frac{a_k}{z^k}.
	\label{joukowski:komplexe-analysis:laurent-reihe}
\end{equation}

% Zusammenhang von |v(z)|^2 und v(z)*^2?

\subsection{Berechnung der Kraft}

Wir können die Gleichung~\eqref{joukowski:komplexe-analysis:laurent-reihe} in den Ausdruck~\eqref{joukowski:auftrieb:auftrieb} einsetzen 
und in ein Integral, mit allen Termen die mit $v_\infty$ multipliziert werden, 
und ein Integral mit allen Termen die nicht mit $v_\infty$ multipliziert werden, umformen, wobei aus

\[
	\begin{aligned}
		F'
		&= 
		-i\frac{\rho}{2} \cdot 
		\displaystyle\oint_{\gamma} v_\infty^2 + 2 v_\infty \sum_{k=1}^{\infty} \frac{a_k}{z^k} + \biggl(\sum_{k=1}^{\infty} \frac{a_k}{z^k}\biggr)^2
		dz \\
		\quad
		&= 
		-i\frac{\rho}{2} \cdot  
		\displaystyle\oint_{\gamma} 
		2 v_\infty \biggl(\frac{v_\infty}{2} + \sum_{k=1}^{\infty} \frac{a_k}{z^k} \biggr)
		dz
		-
		i\frac{\rho}{2} \cdot 
		\displaystyle\oint_{\gamma} 
		\biggl(\sum_{k=1}^{\infty} \frac{a_k}{z^k}\biggr)^2
		dz
	\end{aligned}
\]
wird. 
Beim ersten Ausdruck kann der Term $\frac{v_\infty}{2}$ in $v_\infty - \frac{v_\infty}{2}$ umwandeln. 
Somit lässt sich $-\frac{v_\infty}{2}$ ins zweite Integral verschieben.
Jetzt kann 

\[
	F'
	= 
	-i \rho v_\infty \cdot
	\displaystyle\oint_{\gamma} 
	\biggl(v_\infty + \sum_{k=1}^{\infty} \frac{a_k}{z^k}\biggr)
	dz
	-
	i\frac{\rho}{2} \cdot 
	\displaystyle\oint_{\gamma} 
	\biggl(\sum_{k=1}^{\infty} \frac{a_k}{z^k}\biggr)^2 - \frac{v_\infty^2}{2}
	dz
\]
mit der Gleichung ~\eqref{joukowski:komplexe-analysis:laurent-reihe} in
\[
	F'
	= 
	-i \rho v_\infty \cdot
	\displaystyle\oint_{\gamma} \overline{v(z)}
	dz
	-
	i\frac{\rho}{2} \cdot 
	\displaystyle\oint_{\gamma} 
	\biggl(\sum_{k=1}^{\infty} \frac{a_k}{z^k}\biggr)^2 - \frac{v_\infty^2}{2}
	dz
\]
umgeschrieben werden.
Da $\frac{v_\infty}{2}$ über den ganzen Weg $\gamma$ konstant ist, verschwindet es beim Integrieren um die geschlossene Kurve
und es entsteht
\begin{equation}
	F'
	= 
	-i \rho v_\infty \cdot
	\displaystyle\oint_{\gamma} \overline{v(z)}
	dz
	-
	i\frac{\rho}{2} \cdot 
	\displaystyle\oint_{\gamma} 
	\biggl(\sum_{k=1}^{\infty} \frac{a_k}{z^k}\biggr)^2
	dz.
	\label{joukowski:komplexe-analysis:kraft-mit-zwei-termen}
\end{equation}
Wir definieren für den Ausdruck des ersten Integrals die \emph{Zirkulation} $\Gamma$.

\begin{definition}[Zirkulation]
	\label{joukowski:komplexe-analysis:def:zirkulation}
	Die Zirkulation $\Gamma$ ist das Mass für die um ein Profil drehende Strömung 
	und ist das Ergebnis des Linienintegrals
	\[
		\Gamma
		=
		\displaystyle\oint_{\gamma} \overline{v(z)}
		dz
	\]
	um die geschlossene Kurve $\gamma$.
\end{definition}

Dem zweiten Integral der Gleichung~\eqref{joukowski:komplexe-analysis:kraft-mit-zwei-termen} geben wir den Namen $\Lambda$.
somit kann die Kraft $F'$ vereinfacht geschrieben werden als

\begin{equation}
	F'
	=
	-i \rho v_\infty \Gamma - i \frac{\rho}{2} \Lambda.
	\label{joukowski:komplexe-analysis:kraft-mit-lambda}
\end{equation}

Um die Kraft zu berechnen können wir das Problem in zwei Teile aufteilen:
Einerseits die Auswertung von $\Lambda$,
andererseits die Berechnung der Zirkulation $\Gamma$.

\subsubsection{Berechnung von $\Lambda$}
\label{joukowski:komplexe-analysis:berechnung-von-lambda}

Für die Auswertung von $\Lambda$
aus der Gleichung~\eqref{joukowski:komplexe-analysis:kraft-mit-lambda} können wir ausnutzen,
dass nicht um die Kontour des Flügelprofils direkt integriert werden muss,
sondern dass jeder Weg, welcher das Flügelprofil komplett umschliesst, dasselbe Resultat gibt.
Wir wählen als Weg den Kreis mit Zentrum bei $z_0 = 0$ und mit Radius $r$.
Dafür muss jetzt die Definition~\ref{buch:integration:definition:definition:wegintegral} des Wegintegrals verwendet werden.
Wobei wir die Funktion 
\[
	f(z)
	=
	\biggl(\displaystyle\sum_{k=1}^{\infty} \frac{a_k}{z^k}\biggr)^2
\]
wählen.
Dabei parametrisieren wir $z = \gamma(t) = r e^{i t}$ und stellen das Integral 
\begin{equation}
	\begin{aligned}
		\Lambda
		&=
		\displaystyle\int_{0}^{2\pi} 
		\biggl(\sum_{k=1}^{\infty} \frac{a_k}{r^k e^{k \cdot i t}}\biggr)^2
		i r e^{i t}
		dt \\
		\quad
		&=
		i \displaystyle\int_{0}^{2\pi}
		\bigg(\frac{a_1^2}{r^2 e^{2i t}} + \frac{a_1 a_2}{r^3 e^{3i t}} + \frac{a_2^2}{r^4 e^{4i t}} + \dots \bigg) 
		r e^{i t}
		dt \\
		\quad
		&=
		i \displaystyle\int_{0}^{2\pi}
		\bigg(\frac{a_1^2}{r e^{i t}} + \frac{a_1 a_2}{r^2 e^{2i t}} + \frac{a_2^2}{r^3 e^{3i t}} + \dots \bigg) 
		dt
	\end{aligned}
	\label{joukowski:komplexe-analysis:parametrisiertes-lambda}
\end{equation}
auf.

Da die Auswertung des Integrals nicht vom Radius des Kreises abhängig ist, sofern er das ganze Fügelprofil umschliesst,
können wir den Grenzübergang des Radiuses $\displaystyle\lim_{r\to \infty}$ durchführen.
Weil jeder Term der Gleichung~\eqref{joukowski:komplexe-analysis:parametrisiertes-lambda} 
ein $r$ der Ordnung $1$ bis $n$ im Nenner hat,
gehen alle gegen $0$.

Somit folgt, dass $\Lambda = 0$ sein muss.
Der Term~\eqref{joukowski:komplexe-analysis:kraft-mit-lambda} verkleinert sich zur Form 
\[
	F' = -i \rho v_\infty \Gamma,
\] 
welche als \emph{Satz von Kutta-Joukowski} gilt.

\begin{satz}[Satz von Kutta-Joukowski]
	\label{joukowski:komplexe-analysis:satz:kutta-joukowski}
	Der Satz von Kutta-Joukowski beschreibt die Proportionalität zwischen der komplexen Kraft und der Zirkulation,
	und hat die Form
	\[
		F'
		=
		-i \rho v_\infty \Gamma.
	\]
\end{satz}

\subsubsection{Berechnung der Zirkulation $\Gamma$}
Für die Auswertung von $\Gamma$ kann sehr ähnlich vorgegangen werden, wie bei der Berechnung von $\Lambda$ in Abschnitt~\ref{joukowski:komplexe-analysis:berechnung-von-lambda}.
Auch da wählen wir den Kreis mit Zentrum bei $z_0 = 0$ und mit Radius $r$ und die gleiche parametrisierung von $z$.
Für die Funtion $f(z) = \overline{v(z)}$ setzen wir die Reihenentwicklung aus der Gleichung~\eqref{joukowski:komplexe-analysis:laurent-reihe} ein.
Somit wird aus dem Ausdruck
\[
	\Gamma
	=
	\displaystyle\int_{0}^{2\pi}
	\overline{v(r e^{i t})} \cdot i r e^{i t}
	dt
\]
der Ausdruck
\[
	\begin{aligned}
		\Gamma
		&=
		\displaystyle\int_{0}^{2\pi}
		\biggl(v_\infty + \sum_{k=1}^{\infty} \frac{a_k}{r^k e^{k \cdot i t}}\biggr)
		i r e^{i t}
		dt \\
		\quad
		&=
		i\displaystyle\int_{0}^{2\pi}
		\biggl(v_\infty + \frac{a_1}{r e^{i t}} + \frac{a_2}{r^2 e^{2 i t}} + \frac{a_3}{r^3 e^{3 i t}} + \dots \biggr)
		r e^{i t}
		dt \\
		\quad
		&=
		i\displaystyle\int_{0}^{2\pi}
		\biggl(v_\infty r e^{i t} + a_1 + \frac{a_2}{r e^{i t}} + \frac{a_3}{r^2 e^{2 i t}} + \dots \biggr)
		dt.
	\end{aligned}
\]
Für die Auswertung dieses Integrals gibt es für die einzelnen Terme drei verschiedene Auswertungen.
Der Term $v_\infty r e^{i t}$ integriert gibt
$\biggl[v_\infty r e^{i t}\biggr]_0^{2\pi} = v_\infty r e^{0 \cdot i t} - v_\infty r e^{2\pi \cdot i t}$
wobei $e^{0 \cdot i t } - e^{2\pi \cdot i t} = 0$ ist.
Somit leistet er keinen Beitrag.

Gleich wie bei der Berechnung von $\Lambda$ in Abschnitt~\ref{joukowski:komplexe-analysis:berechnung-von-lambda}
gehen alle Terme $\displaystyle\frac{a_k}{r^{k-1} e^{(k-1) \cdot i t}}$ 
mit $k \in \mathbb{N}$, $k \geq 2$ 
beim Grenzübergang des Radiuses $\displaystyle\lim_{r\to \infty}$ gegen $0$.
Somit leisten auch all diese keinen Beitrag.

Der Einzige Term im Integral, welcher nicht verschwindet ist $a_1$, der zu
\[
	i\displaystyle\int_{0}^{2\pi}
	a_1
	dt
	=
	a_1 2\pi i
\]
ausgewertet wird.
Somit gilt
\begin{equation}
	\Gamma
	=
	a_1 2\pi i.
	\label{joukowski:komplexe-analysis:zirkulation-berechnet}
\end{equation}
Mit diesem Resultat können wir endlich die komplexe Kraft, welche auf das Flügelprofil wirk ausrechnen.
Durch das Einsetzen von Gleichung~\eqref{joukowski:komplexe-analysis:zirkulation-berechnet}
in Satz~\ref{joukowski:komplexe-analysis:satz:kutta-joukowski}
entsteht der Term
\begin{equation}
	F'
	=
	-i \rho v_\infty \cdot a_1 2\pi i
	= \rho v_\infty a_1 2 \pi.
	\label{joukowski:komplexe-analysis:auftriebskraft-berechnet}
\end{equation}
% Ist a_1 elem C, sonst keine auftriebskraft?